\chapter{Glosario de términos}

Este anexo recoge los principales términos técnicos utilizados a lo largo de la memoria, con el objetivo de facilitar la lectura del documento y unificar el significado de los conceptos más relevantes del proyecto.

\begin{longtable}{p{0.25\textwidth}p{0.68\textwidth}}
\caption{Glosario de términos del proyecto}\label{tab:glosario}\\
\toprule
\textbf{Término} & \textbf{Definición} \\
\midrule
\endfirsthead

\toprule
\textbf{Término} & \textbf{Definición} \\
\midrule
\endhead

\midrule
\multicolumn{2}{r}{Continúa en la página siguiente} \\
\endfoot

\bottomrule
\endlastfoot

\textbf{API REST} & Interfaz de comunicación basada en HTTP que permite al frontend consultar y modificar información gestionada por el backend. \\

\textbf{Backend} & Parte del sistema encargada de procesar la lógica de negocio, exponer la API, validar peticiones y comunicarse con la base de datos y servicios externos. \\

\textbf{Cloudinary} & Servicio externo utilizado para almacenar y servir imágenes asociadas a eventos, chats y perfiles de usuario. \\

\textbf{DTO} & Objeto de transferencia de datos utilizado en el backend para definir y validar la estructura de la información recibida desde el cliente. \\

\textbf{Expo} & Plataforma y conjunto de herramientas que facilitan el desarrollo, prueba y distribución de aplicaciones React Native. \\

\textbf{Firebase Authentication} & Servicio empleado para gestionar el registro, inicio de sesión y verificación de identidad de los usuarios. \\

\textbf{Frontend} & Parte cliente de la aplicación, desarrollada como aplicación móvil, con la que interactúa directamente el usuario. \\

\textbf{Geocodificación} & Proceso de convertir una dirección textual en coordenadas geográficas que puedan almacenarse y consultarse en el sistema. \\

\textbf{GitHub Actions} & Plataforma de integración continua utilizada para automatizar comprobaciones de calidad y flujos de despliegue. \\

\textbf{Mockup} & Representación visual preliminar de una pantalla o flujo de interfaz, utilizada durante la fase de análisis y diseño. \\

\textbf{NestJS} & Framework de backend para Node.js utilizado para construir la API REST del sistema. \\

\textbf{OpenStreetMap} & Proyecto cartográfico abierto empleado como base para mapas y servicios de geolocalización. \\

\textbf{PostGIS} & Extensión de PostgreSQL que permite almacenar y consultar datos geográficos de forma eficiente. \\

\textbf{React Native} & Framework utilizado para desarrollar la aplicación móvil multiplataforma a partir de una base de código común. \\

\textbf{Scraping} & Técnica utilizada para extraer eventos desde páginas web o fuentes externas y transformarlos en datos estructurados. \\

\textbf{Sprint} & Iteración de trabajo dentro de la metodología ágil empleada para planificar, desarrollar y revisar funcionalidades. \\

\textbf{TypeORM} & Herramienta ORM utilizada para mapear entidades TypeScript con tablas de la base de datos PostgreSQL. \\

\textbf{TypeScript} & Lenguaje de programación empleado tanto en frontend como en backend, que añade tipado estático sobre JavaScript. \\

\textbf{WebSocket} & Tecnología de comunicación bidireccional utilizada para funcionalidades en tiempo real, como el chat. \\
\end{longtable}
